\documentclass[12pt]{article}
\usepackage{../../../lecture_notes}
\usepackage{../../../math}
\usepackage{../../../uark_colors}

\hypersetup{
  colorlinks = true,
  allcolors = ozark_mountains,
  breaklinks = true,
  bookmarksopen = true
}

\begin{document}
\begin{center}
  {\Huge\bf Quiz --- Assignment 2}

  \smallskip
  {\large\it ECON 5783 --- University of Arkansas}
\end{center}

\vspace{5mm}
\noindent\textbf{Name:} \rule{0.7\textwidth}{0.4pt}

\bigskip
\begin{enumerate}
  \item Let
  $$
    m(t) \equiv \expec{\text{Exam Score}_i}{\text{Tutoring Hours}_i = t}.
  $$
  In words, what does $m(4)$ represent?
  Describe how you could estimate it from a sample without imposing a linear relationship between tutoring hours and exam scores.

  
  \vspace*{10\bigskipamount}
  \item An analyst models exam scores as a linear function of tutoring hours:
  $$
    \text{Exam Score}_i = \alpha + \beta \text{Tutoring Hours}_i + u_i.
  $$
  How does this model restrict the conditional expectation function compared with a fully flexible model?

  
%   \vspace*{\bigskipamount}
%   \item \pts{4} The following regression relates monthly electricity use to home type and electric-vehicle ownership.
%   Assume \texttt{home\_type} includes \texttt{Apartment}, \texttt{Detached}, and \texttt{Townhouse}.
%   What indicators does \texttt{i()} create, which home type is omitted, and how can the relationship between EV ownership and electricity use vary across home types?
%   \begin{codeblock}
% feols(monthly_kwh ~ i(home_type, ref = "Apartment"), data = homes)
%   \end{codeblock}
\end{enumerate}

\end{document}
